%% ch18 --- Po co ci chaos: przewodnik terenowy po przewidywaniu
%%
%% Napisany we wrześniu 2026 jako rozdział zamykający książkę. Zamienia
%% książkę w pięć pytań, które warto zadać każdemu przewidywaniu (i dwa
%% kolejne, zebrane w przewodniku), a na każde odpowiada liczbą policzoną
%% TUTAJ: każda wartość pochodzi z code/measure/ch18.py, który importuje TE
%% SAME funkcje, które drukują listingi (code/ch18/).
%%
%% CZEGO TEN ROZDZIAŁ NIE MÓWI. Nie cytuje żadnej liczby zmierzonej w innym
%% rozdziale. Wykładniki, horyzonty i szanse są tu liczone od nowa, a liczba,
%% której książka nie mierzy (dwa tygodnie pogody, pięciomilionowy czas
%% Lapunowa planet, sekundy podwójnego wahadła), jest oznaczona jako
%% cytowana albo jako rząd wielkości.
%%
%% Bliźniak: chapters/en/ch18-why.tex. Podrozdział za podrozdziałem, makro
%% za makrem, liczba za liczbą; tools/parity.py je porównuje.

\chapter{Po co ci chaos: przewodnik terenowy po przewidywaniu}
\label{ch:why}
\index{przewidywanie!przewodnik terenowy}\index{prognoza}

\noindent
Oto cztery przewidywania, na które możesz trafić w ciągu jednego tygodnia.
Tablica mówi, że całkowite zaćmienie Słońca nastąpi 2 sierpnia 2027 roku.
Aplikacja pogodowa podaje 70 procent szans na deszcz w sobotę. Wyciąg z
funduszu emerytalnego mówi, ile będą warte twoje oszczędności w 2060 roku.
A program komputerowy odpowiada na pytanie o przyszłoroczne stopy
procentowe bez cienia wątpliwości.

Któremu z nich zaufasz i jak daleko? Teraz możesz na to odpowiedzieć
liczbami, które sam policzyłeś. Ten rozdział zamienia książkę w przewodnik
terenowy: pięć pytań, które warto zadać każdemu przewidywaniu, każde z
odpowiedzią opartą na pomiarze, a potem argument, dlaczego każdy ich
potrzebuje.

\begin{outcomes}
\outcome{zapytać o każde przewidywanie, czy mała zmiana stanu początkowego
  rośnie}
\outcome{policzyć horyzont prognozy i powiedzieć, co dają lepsze dane}
\outcome{powiedzieć, co da się przewidzieć za horyzontem, i czytać szanse}
\outcome{odróżnić chaos od szumu, a małą przyczynę od dużego skutku}
\outcome{zastosować siedem pytań przewodnika do prognozy albo modelu}
\outcome{wyjaśnić, liczbami, które sam policzysz, dlaczego jedne rzeczy da
  się przewidzieć na stulecia, a inne tylko na dni}
\end{outcomes}

\section{Pytanie pierwsze: czy mała zmiana rośnie?}
\label{sec:ch18-grow}
\index{wrażliwość na warunki początkowe!test}

\noindent
Zacznij od najczystszej pary z listy: zaćmienia i deszczu. Obie prognozy
powstają tak samo. Komputer bierze zmierzoną teraźniejszość, stosuje
dokładne prawa fizyki i idzie krok po kroku naprzód w czasie.

\begin{think}
Tablica zaćmień jest dobra na tysiące lat, a szczegółowa prognoza pogody na
mniej więcej tydzień. Obie powstają z dokładnych praw, stosowanych krok po
kroku od zmierzonego stanu początkowego. Na czym polega różnica, w języku tej
książki?
\end{think}
\reveal{Na tym, jak szybko rośnie mały błąd stanu początkowego. W ruchach
Ziemi i Księżyca rośnie tak wolno, że przez tysiąclecia pozostaje mały; w
powietrzu rośnie, aż staje się tak duży jak sama pogoda.}
Sprawdza się to testem, który rozdział~\ref{ch:butterfly} zamienił w pomiar:
uruchom regułę dwa razy, z dwóch stanów początkowych bliższych sobie, niż
potrafiłby je rozróżnić jakikolwiek przyrząd, i patrz na różnicę. Oto trzy
reguły, które już znasz z tej książki, każda uruchomiona dwa razy z
odległości jednej miliardowej:

\pyregion{code/ch18/grows.py}{twice}{Pierwsze pytanie, zadane trzem regułom}{lst:ch18-twice}

\transcript{ch18-grows}

\noindent
Stygnąca herbata zapomina o zaburzeniu: różnica jest co minutę mnożona
przez \num{0.9}. Odwzorowanie logistyczne przy $r = \num{2.8}$ też
zapomina, bo osiada w swoim punkcie stałym. Przy $r = 4$ jedna miliardowa
rośnie, aż staje się tak duża, jak pozwala samo odwzorowanie. Jeśli
różnica maleje, trajektorię da się przewidzieć tak długo, jak długo reguła
jest poprawna. Jeśli rośnie, następne pytanie brzmi: jak szybko?

Pytanie pierwsze ma przed sobą jeszcze jedno: czy w ogóle istnieje dokładna
reguła? Dla planet i dla trzech równań Lorenza istnieje. W gospodarce czy w
epidemii sama reguła zmienia się, gdy ludzie reagują na prognozę, a to jest
inna trudność niż chaos.

\section{Pytanie drugie: jak daleko jest horyzont?}
\label{sec:ch18-horizon}
\index{horyzont prognozy}\index{wykładnik Lapunowa}

\noindent
Rosnący błąd jest zwykle w każdym kroku mnożony przez mniej więcej ten sam
czynnik, a rozdział~\ref{ch:lyapunov} zamienił ten czynnik w jedną liczbę: wykładnik
Lapunowa, zapisywany $\lambda$ (grecka litera lambda), czyli logarytm
naturalny $\ln$ średniego czynnika. Czynnik $2$ to wykładnik $\ln 2$, mniej
więcej \val{ch18.ln.two}. Załóżmy, że stan początkowy jest błędny o $e_0$,
a prognoza staje się bezużyteczna, gdy błąd przekroczy tolerancję $T$. Po
$n$ krokach błąd to $e_0$ pomnożone $n$ razy przez ten czynnik, więc sięga
$T$, gdy
\[ n \approx \frac{\ln(T / e_0)}{\lambda}. \]
Słowami: horyzont to logarytm tego, ile razy za mały jest twój błąd,
podzielony przez wykładnik. Biblioteka książki ma to w jednym wierszu
(\code{e**exponent} to tam czynnik, odtworzony z wykładnika):

\pyregion{code/src/chaoslab/measure.py}{horizon}{Wzór na horyzont w bibliotece książki}{lst:ch18-horizon-lib}

\noindent
Wypróbuj go na odwzorowaniu logistycznym przy $r = 4$, którego wykładnik ten
rozdział mierzy jako \val{ch18.lam.logistic} na krok ($\ln 2$, jak pokazuje
rozdział~\ref{ch:lyapunov}). Przyjmij, że znasz stan początkowy z dokładnością do
sześciu miejsc po przecinku, więc $e_0$ to jedna milionowa, i uznaj prognozę za
bezużyteczną, gdy myli się o więcej niż $T = \num{0.1}$. Wtedy
$T / e_0 = 100\,000$, a $n \approx \val{ch18.ln.six} /
\val{ch18.lam.logistic} \approx \val{ch18.hor.six}$ kroku.

\begin{think}
Kupujesz przyrząd tysiąc razy dokładniejszy i teraz znasz stan początkowy z
dokładnością do dziewięciu miejsc po przecinku. Jak długo wytrzyma
prognoza: tysiąc razy dłużej, mniej więcej trzy razy dłużej, czy o kilka
kroków dłużej?
\end{think}
\reveal{O kilka kroków dłużej: mniej więcej o \val{ch18.hor.gain} więcej,
czyli razem \val{ch18.hor.nine}. Podzielenie $e_0$ przez tysiąc dodaje
$\ln 1000 / \lambda = \val{ch18.ln.thousand} / \val{ch18.lam.logistic}$
kroku i nic więcej.}
To odpowiedź wzoru. Oto odpowiedź samego odwzorowania, zmierzona na
piechotę dla \val{ch18.apart.starts} różnych stanów początkowych:

\pyregion{code/ch18/thousandfold.py}{measure}{Tysiąc razy lepsze dane, zmierzone}{lst:ch18-thousandfold}

\transcript{ch18-thousandfold}

\plotfig{ch18-digits}{Jak daleko sięga prognoza w zależności od tego, ile
miejsc po przecinku stanu początkowego jest poprawnych. Linie to wzór,
punkty to mediana wielu zmierzonych par przebiegów. Każda dodatkowa cyfra
kupuje tych samych kilka kroków.}{horyzont a liczba poprawnych cyfr}

\noindent
Wzór mówił \val{ch18.hor.six} i \val{ch18.hor.nine}; odwzorowanie mówi
\val{ch18.apart.six} i \val{ch18.apart.nine}, czyli o
\val{ch18.apart.gain} więcej. Każda dodatkowa poprawna
cyfra kupuje tę samą liczbę kroków, $\ln 10 / \lambda = \val{ch18.ln.ten} /
\val{ch18.lam.logistic}$, tutaj mniej więcej \val{ch18.per.digit}. Horyzont
rośnie jak logarytm twojej dokładności, więc lepsze przyrządy i większe
komputery mogą tylko dodać stały odcinek. Mały kalkulator zestawia obok
siebie trzy układy z tej książki:

\pyregion{code/ch18/horizon.py}{calculator}{Kalkulator horyzontu}{lst:ch18-calculator}

\transcript{ch18-horizon}

\noindent
Błędy odwzorowania Hénona rosną wolniej, z wykładnikiem \val{ch18.lam.henon}
na krok, więc każda cyfra kupuje mu
\val{ch18.per.digit.henon} kroku; układ Lorenza, przy
\val{ch18.lam.lorenz} na jednostkę czasu, zyskuje \val{ch18.per.digit.lorenz}
jednostki czasu na cyfrę. Na wykresie w zależności od liczby poprawnych
cyfr każdy horyzont jest linią prostą (rysunek~\ref{fig:ch18-digits}).

Ta sama arytmetyka wyjaśnia zaćmienie. Obliczenia Układu Słonecznego, które
Jacques Laskar przeprowadził w 1989 roku, przytoczone w
rozdziale~\ref{ch:mechanics}, dały planetom wewnętrznym wykładnik około
jednej piątej na milion lat: czas Lapunowa, czyli odwrotność wykładnika,
wynosi mniej więcej pięć milionów lat. Przyjmij, że znasz położenie Ziemi z
dokładnością do metra, i uznaj prognozę za bezużyteczną, gdy błąd dorówna rozmiarom jej
orbity, 150 milionom kilometrów. Horyzont to pięć milionów lat razy
$\ln 150\,000\,000\,000 \approx \val{ch18.planets.ln}$, czyli jakieś
\val{ch18.planets.hor} milionów lat, a milimetr dodałby tylko około
\val{ch18.planets.gain} milionów. Układ Słoneczny jest chaotyczny; w skali
tablicy zaćmień jego błędy jeszcze nie zaczęły mieć znaczenia.

\begin{lab}{l18_01_horizon}{Kalkulator horyzontu}
Napisz \code{forecast\_length}, która zamienia wykładnik i liczbę poprawnych
cyfr w horyzont; \code{extra\_steps}, która mówi, co daje dokładniejszy
pomiar; oraz \code{digits\_needed}, która stosuje wzór w drugą stronę.
Potem zapytaj, ilu cyfr potrzebowałoby sto kroków odwzorowania
logistycznego, i porównaj to z mniej więcej szesnastoma cyframi dziesiętnymi,
które mieszczą zwykłe liczby komputera (rozdział~\ref{ch:computers}).
\end{lab}

\section{Pytanie trzecie: co przetrwa horyzont?}
\label{sec:ch18-beyond}
\index{prognoza zespołowa}\index{atraktor!statystyki}

\noindent
Za horyzontem trajektoria jest stracona. To nie to samo co brak wszelkiej
wiedzy, i właśnie tę część teorii chaosu najczęściej się przeocza.

\begin{think}
Czterdzieści kroków naprzód trajektoria odwzorowania logistycznego jest
dawno stracona. Czy o kroku czterdziestym da się jeszcze coś
przewidzieć?
\end{think}
\reveal{Tak: szanse. Szansa, że $x$ jest poniżej \num{0.25} w kroku
czterdziestym, wynosi mniej więcej jedną trzecią, i taki sam jest ułamek
czasu, który spędza tam każdy długi przebieg. Trajektoria przepadła; jej
statystyki nie.}
Twórcy prognoz uzyskują takie szanse, uruchamiając model wiele razy, ze stanów
początkowych rozłożonych po wszystkim, na co pozwala pomiar, i licząc. To
jest zespół, a rozdział~\ref{ch:weather} buduje go dla pogody. Oto
najmniejsza wersja:

\pyregion{code/ch18/beyond.py}{ensemble}{Prognoza z zespołu i jeden długi przebieg}{lst:ch18-ensemble}

\transcript{ch18-beyond}

\noindent
Gdy stan początkowy jest poprawny do sześciu cyfr, wszystkie przebiegi,
a jest ich \val{ch18.ens.members}, zgadzają się aż do kroku
\val{ch18.ens.sure.six}, a prognoza jest jedną pewną odpowiedzią. Potem się
rozchodzą, prognoza staje się szansami, a po kilku krokach szanse ustalają
się na \val{ch18.ens.clim}: ułamku czasu, który jeden długi przebieg spędza
poniżej \num{0.25}, z każdego z dwóch stanów początkowych. Przy dziewięciu
cyfrach przebiegi zgadzają się aż do kroku \val{ch18.ens.sure.nine}, a
potem dzieje się to samo (rysunek~\ref{fig:ch18-odds}).

\plotfig{ch18-odds}{Prognoza zespołowa szansy, że $x$ jest poniżej
\num{0.25}. Jest pewna, dopóki przebiegi się zgadzają, podaje szanse, gdy
się rozchodzą, i ustala się na ułamku z długiego przebiegu. Trzy cyfry
więcej to opóźniają, ale nie powstrzymują.}{prognoza, która staje się szansami}

Prognoza układu chaotycznego ma więc trzy postacie, jedna po drugiej:
trajektorię, dopóki przebiegi się zgadzają, szanse, gdy się rozchodzą, i
potem statystyki długiego przebiegu. Te ostatnie rozdział~\ref{ch:lorenz}
nazywa statystykami atraktora, a w przypadku pogody to jest klimat.

\begin{trapbox}
\textbf{\enquote{Nieprzewidywalne znaczy niepoznawalne.}} Za horyzontem nie
da się poznać trajektorii, a prawie wszystko inne tak: dokąd układ może
dojść, a dokąd nie, jak często odwiedza każdy obszar, jakie są szanse w
każdym pytaniu, jakie zechcesz zadać. 70 procent szans na deszcz to nie
wykręt synoptyka; to poprawna odpowiedź na pytanie, na które chaos
uniemożliwił odpowiedź \enquote{tak} albo \enquote{nie}. Projekcja
klimatyczna prognozuje statystyki, nie wtorki.
\end{trapbox}

\begin{worldbox}
\enquote{Stożek niepewności}, który amerykańskie Narodowe Centrum
Huraganów (National Hurricane Center) rysuje wokół prognozowanego toru
burzy, rozszerza się z każdym dniem wyprzedzenia. Jego rozmiar wyznaczają
błędy niedawnych prognoz samego centrum, tak aby środek burzy pozostawał w
nim mniej więcej dwa razy na trzy: horyzont i zestaw szans na jednym
obrazku, publikowany po to, by działały według niego miliony ludzi.
\end{worldbox}

\section{Pytanie czwarte: chaos, szum czy mała przyczyna?}
\label{sec:ch18-noise}
\index{szum}\index{efekt motyla}

\noindent
Nieregularne dane zachęcają do opowieści, a dwie z nich wymagają osobnego
pytania.

Pierwsza mówi: \enquote{to jest chaos}. Rozdział~\ref{ch:life} pokazuje, że
szereg chaotyczny i losowy mogą mieć ten sam histogram i że rozróżnia je
struktura: nanieś każdą wartość względem następnej, a chaos narysuje
krzywą tam, gdzie szum rysuje chmurę, i chaos da się prognozować krok lub
dwa naprzód tam, gdzie szumu nie. Ma znaczenie, co masz, i ile jest danych:
krótki, zaszumiony zapis może wyglądać prawie jak cokolwiek.

Druga opowieść biegnie wstecz, od dużego skutku do małej przyczyny.

\begin{think}
Pchnij stygnącą herbatę i odwzorowanie logistyczne przy $r = 4$ o jedną
miliardową każde, jak w podrozdziale~\ref{sec:ch18-grow}. Które pchnięcie
coś zmieniło czterdzieści kroków później? Czy zmieniło ułamek czasu, który
odwzorowanie spędza poniżej \num{0.25}?
\end{think}
\reveal{Tylko pchnięcie odwzorowania logistycznego i tylko jego trajektorię:
różnica urosła do rozmiarów samego odwzorowania. Ułamek czasu poniżej
\num{0.25} to wciąż mniej więcej jedna trzecia, z każdego stanu
początkowego. Pchnięcie herbaty wygasło.}

\begin{trapbox}
\textbf{\enquote{Małe przyczyny zawsze mają znaczenie.}} W układzie
stabilnym wygasają, jak w przypadku herbaty. W chaotycznym mała przyczyna
zmienia szczegóły za horyzontem i nic w statystykach, a ponieważ każda inna
mała przyczyna też zmienia szczegóły, żadnej nie da się potem wskazać jako
powodu. Motyl z rozdziału~\ref{ch:butterfly} nie spowodował burzy.
Obwinianie jednej małej przyczyny o duży skutek po fakcie to opowieść, nie
pomiar.
\end{trapbox}

\section{Pytanie piąte: czy da się tym sterować?}
\label{sec:ch18-steer}
\index{sterowanie chaosem}\index{synchronizacja}

\noindent
Ostatnie pytanie zamienia ograniczenie w narzędzie.

\begin{think}
Jeśli drobna zmiana stanu początkowego całkowicie zmienia przyszłość, to
czy jest to zła, czy dobra wiadomość dla kogoś, kto chce układem sterować?
\end{think}
\reveal{Jedno i drugie. Właśnie dlatego nie da się prognozować układu
daleko naprzód i właśnie dlatego drobne, dobrze wymierzone pchnięcie może go
przesunąć daleko.}
Edward Ott, Celso Grebogi i James Yorke uczynili z tego ideę swojej pracy o
sterowaniu chaosem z 1990 roku: poczekaj, aż układ przejdzie blisko
jednego z niestabilnych cykli ukrytych w jego atraktorze, lekko popchnij
parametr, a da się go utrzymać na tym cyklu. Rozdział~\ref{ch:taming} robi
to pchnięciami mniejszymi niż jeden procent i sprawia, że dwa chaotyczne
układy idą w równym kroku, jak pokazali w tym samym roku Louis Pecora i
Thomas Carroll. Chaos nie jest więc wymówką dla fatalizmu. Układ, którego
nie da się prognozować, może być układem, którym da się sterować, a zmiana
parametru, na przykład wymuszenia w modelu klimatu, przesuwa statystyki
nawet wtedy, gdy żadnej trajektorii nie da się przewidzieć.

\section{Przewodnik terenowy na jednej stronie}
\label{sec:ch18-guide}
\index{przewidywanie!przewodnik terenowy}

\noindent
Oto pięć pytań, dwa kolejne, na które książka zapracowała, i miejsce, w
którym zmierzyłeś każde z nich:

\begin{enumerate}[itemsep=2pt]
\item Czy istnieje dokładna reguła i czy mała zmiana stanu początkowego
  rośnie? (rozdział~\ref{ch:butterfly})
\item Jak daleko jest horyzont i co dałyby lepsze dane?
  (rozdział~\ref{ch:lyapunov})
\item Co przetrwa horyzont: atraktor, statystyki, szanse?
  (rozdziały~\ref{ch:lorenz} i~\ref{ch:weather})
\item Czy nieregularność to chaos, czy szum, i czy danych starczy, żeby to
  rozstrzygnąć? (rozdział~\ref{ch:life})
\item Czy da się tym sterować albo przesunąć statystyki parametrem?
  (rozdział~\ref{ch:taming})
\item Czy odpowiedź przetrwa mniejszy krok i więcej cyfr?
  (rozdziały~\ref{ch:motion} i~\ref{ch:computers})
\item Czy ktoś po fakcie obwinia małą przyczynę o duży skutek?
  (rozdział~\ref{ch:butterfly})
\end{enumerate}

\mermaidfig{ch18-field-guide}{Przewodnik terenowy w trzech ruchach. Jeśli
mały błąd nie rośnie, przetrwa sama trajektoria; jeśli rośnie, horyzont
mówi, jak długo, a za nim pytasz o statystyki i szanse.}{przewodnik
terenowy w trzech ruchach}

\noindent
A oto znajome prognozy uporządkowane według horyzontu. Tylko jedna liczba
jest policzona w tym rozdziale; reszta to rzędy wielkości, niektóre z nich
z obiegowej wiedzy, oznaczone według tego, skąd pochodzą.

\medskip
\noindent
{\small
\begin{tabularx}{\linewidth}{@{}>{\raggedright\arraybackslash}X
  >{\raggedright\arraybackslash}p{2.9cm}
  >{\raggedright\arraybackslash}p{2.9cm}
  >{\raggedright\arraybackslash}X@{}}
\toprule
\textbf{Prognoza} & \textbf{Horyzont} & \textbf{Skąd ta liczba}
  & \textbf{Za nim pytaj o} \\
\midrule
Gdzie są planety & około \val{ch18.planets.hor} milionów lat
  & policzone wyżej & gdzie leżą ich orbity \\ \addlinespace[3pt]
Zaćmienia & tysiące lat & opublikowane tablice & nic więcej \\ \addlinespace[3pt]
Pływy astronomiczne & lata & opublikowane tablice
  & falę sztormową, która jest pogodą \\ \addlinespace[3pt]
Szczegółowa pogoda & tydzień w praktyce, najwyżej dwa tygodnie
  & cytowane (Lorenz, 1969) & szanse, potem klimat \\ \addlinespace[3pt]
Podwójne wahadło & sekundy & rząd wielkości
  & dokąd pozwala mu jego energia \\ \addlinespace[3pt]
Kurs akcji z dnia na dzień & nieustalony & nieregularny to nie chaotyczny
  & rozrzut dawnych zmian \\ \addlinespace[3pt]
\bottomrule
\end{tabularx}\par}

\medskip
\noindent
Teraz wyciąg emerytalny i pewny siebie program. Projekcja emerytalna to
wynik, jaki dałaby założona stopa zwrotu, więc pytaj o zakres stóp. Pewna
siebie odpowiedź o przyszłorocznej gospodarce nie ma horyzontu, dopóki ktoś
nie powie, jaka jest jej reguła, a pewny ton to nie przedział błędu.

\begin{lab}{l18_02_guide}{Przewodnik terenowy na modelu, który sam piszesz}
Napisz \code{error\_grows}, \code{growth\_rate} i
\code{long\_run\_fraction}: trzy pierwsze pytania jako kod, który możesz
skierować na dowolną regułę zamieniającą jedną liczbę w następną. Test
sprawdza je na stygnącej herbacie i odwzorowaniu logistycznym. Potem
napisz na końcu pliku własną regułę (namiot, sinus, populację z odłowem) i
zadaj jej wszystkie siedem pytań.
\end{lab}

\section{Dlaczego każdy tego potrzebuje}
\label{sec:ch18-everyone}
\index{chaos!dlaczego ma znaczenie}

\noindent
\emph{Chaos zakończył erę mechanizmu zegara}, nie dodając ani jednego losowego
składnika: dokładna reguła i prawie dokładny stan początkowy dają przyszłość
znaną tylko do horyzontu, który rośnie jak logarytm twojej dokładności.

\emph{Mówi, którym prognozom ufać i jak długo.} Zaćmienie i burza podlegają
prawom tego samego rodzaju; różni je wykładnik. To uczy pokory wobec
pewności na długi dystans, bo pojedyncza liczba daleko za horyzontem to
zgadywanie przebrane za pomiar, i zaufania do modeli na krótki dystans, bo
w obrębie horyzontu są zdumiewająco dobre.

\emph{Czyni prawdopodobieństwa szacownymi.} Szanse to postać, jaką za
horyzontem przybiera poprawna prognoza, a nie porażka prognozowania.

\emph{Jest piękny i uniwersalny.} Stała Feigenbauma wychodzi z linijki
algebry i z prawdziwych cieczy i obwodów (rozdział~\ref{ch:bifurcations});
fraktale nadają linii brzegowej wymiar między jeden a dwa
(rozdział~\ref{ch:fractals}); zbiór Mandelbrota mieści całą drogę do
chaosu (rozdział~\ref{ch:mandelbrot}).

\emph{I jest zasobem}: sterowanie, synchronizacja i wymieszanie, jakie mają
mieć liczby losowe komputera (choć rozdział~\ref{ch:computers} ostrzega, że
samo odwzorowanie chaotyczne jest ich kiepskim źródłem).

Teoria chaosu nie jest teorią wszystkiego ani nieładu w ogóle: większość
nieregularności świata to szum, niewiedza albo reguła, która ciągle się
zmienia.

\begin{think}
Kolega mówi: \enquote{Gospodarka jest chaotyczna, więc każda prognoza jest
równie dobra jak inna}. Co jest w tym źle?
\end{think}
\reveal{Obie połowy. Nikt nie wykazał, że gospodarka podlega stałej
chaotycznej regule, a nawet układ chaotyczny nie może robić czegokolwiek:
trzyma się swojego atraktora, jego szanse są ustalone, a jego prognoza na
krótki dystans może być bardzo dobra.}

\begin{trapbox}
\textbf{\enquote{Chaos znaczy: wszystko wolno.}} Chaos to dokładna reguła,
której ruch pozostaje w ograniczonym atraktorze, ze stałymi statystykami i
horyzontem, który da się policzyć. Nazwanie czegoś chaotycznym to
twierdzenie, które da się sprawdzić, a nie licencja na przyklejanie
motylich skrzydeł wszystkiemu, co nieprzewidywalne, i nie powód, żeby
przestać prognozować.
\end{trapbox}

\section{Demon Laplace'a po raz ostatni}
\label{sec:ch18-demon}
\index{demon Laplace'a}

\noindent
Rozdział~\ref{ch:models} zakończył się zasianiem pytania, a teraz możesz
na nie odpowiedzieć.

\begin{think}
Załóżmy, że reguła jest dokładna i znasz ją doskonale. Czy przyszłość jest
wtedy znana, dowolnie daleko naprzód?
\end{think}
\reveal{Nie. Jest znana tylko do horyzontu, który rośnie jak logarytm tego,
jak dokładnie znasz stan początkowy. Za nim znany jest atraktor, statystyki
i szanse.}
Demon Laplace'a potrzebuje nieskończenie dokładnej teraźniejszości, a każda
dodatkowa cyfra kupuje mu tylko kilka kroków więcej. Mechanizm zegara jest
prawdziwy dla zaćmienia i pływów, a kończy się dla pogody, podwójnego
wahadła i jednej linijki algebry.

Zostają prognozy, które spotykasz codziennie: pociąg za siedem minut,
deszcz w sobotę, projekcja, pewna siebie odpowiedź. Zadaj każdej pięć
pytań. Niektórym prognozom zaufasz bardziej niż dotąd, innym znacznie
mniej, i po raz pierwszy będziesz umieć powiedzieć dlaczego, liczbą.

\begin{summarybox}
\begin{itemize}
\item \textbf{Pytanie pierwsze}: uruchom regułę dwa razy ze stanów
  początkowych odległych o włos. Jeśli różnica maleje, trajektorię da się
  prognozować tak daleko, jak długo reguła jest poprawna.
\item \textbf{Pytanie drugie}: \textbf{horyzont} to mniej więcej
  $\ln(T/e_0)$ podzielone przez wykładnik $\lambda$, więc każda poprawna
  cyfra kupuje tych samych kilka kroków.
\item \textbf{Pytanie trzecie}: za horyzontem przetrwają
  \textbf{statystyki} i \textbf{szanse}, a zespół pokazuje, jak trajektoria
  zamienia się w szanse.
\item \textbf{Pytanie czwarte}: nieregularny to nie chaotyczny, dopóki nie
  pokaże tego struktura, a dużego skutku nie da się potem przypisać jednej
  małej przyczynie.
\item \textbf{Pytanie piąte}: wrażliwość czyni chaos
  \textbf{sterowalnym}.
\item \textbf{Przewodnik terenowy} dodaje dwa pytania i porządkuje znajome
  prognozy według horyzontu.
\item Chaos zakończył erę \textbf{mechanizmu zegara} i czyni szanse
  szacownymi; nie jest licencją na \enquote{wszystko wolno}.
\item \textbf{Demon Laplace'a} znałby przyszłość tylko do horyzontu.
\end{itemize}
\end{summarybox}

\canyou

\begin{checkpoint}
\item Dwa przebiegi reguły zaczynają się w odległości jednej miliardowej;
  trzydzieści kroków później dzieli je jedna bilionowa. Co mówi pytanie
  pierwsze?
  \answerto{Błąd maleje, więc trajektorię da się prognozować tak daleko, jak
  długo sama reguła jest poprawna.}
\item Błędy w pewnym układzie rosną z wykładnikiem \num{0.5} na dzień. Znasz
  stan początkowy z dokładnością do \num{0.001} i możesz tolerować błąd $1$.
  Wiedząc, że $\ln 1000 \approx \val{ch18.ln.thousand}$, podaj horyzont.
  \answerto{Mniej więcej $\val{ch18.ln.thousand} / \num{0.5} \approx
  \num{13.8}$ dnia.}
\item Teraz mierzysz ten stan początkowy dziesięć razy dokładniej. Wiedząc,
  że $\ln 10 \approx \val{ch18.ln.ten}$, powiedz, o ile dłużej wytrzyma
  prognoza.
  \answerto{Mniej więcej o $\val{ch18.ln.ten} / \num{0.5} \approx \num{4.6}$
  dnia dłużej: stały dodatkowy odcinek, nie dziesięć razy dłużej.}
\item Daleko za horyzontem zespół przebiegów odwzorowania logistycznego przy
  $r = 4$ daje szansę jedną na trzy, że $x$ jest poniżej \num{0.25}. Czy
  ta prognoza jest bezwartościowa?
  \answerto{Nie. To poprawna prognoza: ułamek czasu, który odwzorowanie tam
  spędza, taki sam z każdego stanu początkowego.}
\item Szereg codziennych pomiarów wygląda nieregularnie. Wymień dwie rzeczy,
  które trzeba sprawdzić, zanim nazwie się go chaosem.
  \answerto{Czy każda wartość naniesiona względem następnej rysuje krzywą,
  a nie chmurę; czy da się go prognozować krok lub dwa naprzód z podobnych
  chwil z przeszłości; i czy zapis jest dość długi, by to rozstrzygnąć.}
\item Dlaczego ta sama własność, która rujnuje dalekie prognozy układu
  chaotycznego, pozwala nim sterować?
  \answerto{Układ jest wrażliwy na małe zmiany: drobny błąd rośnie, i tak
  samo rośnie skutek drobnego, dobrze wymierzonego pchnięcia.}
\item Wyjaśnij, dlaczego zaćmienia da się przewidzieć na tysiące lat, a
  pogodę tylko na mniej więcej tydzień.
  \answerto{Horyzont to $\ln(T/e_0)$ podzielone przez wykładnik. Wykładnik
  Układu Słonecznego to mniej więcej jedna piąta na milion lat, co daje
  horyzont około stu milionów lat; błędy atmosfery rosną tak samo w ciągu
  dni. Lepsze dane dodają do każdego z nich stały odcinek.}
\end{checkpoint}
